% Vorlage für WABs der Provadis Hochschule
% basiert auf "Universität Ulm Praktikumsbericht Vorlage"
% von Max Sch.
% CC BY 4.0
%
% Einstellungen.tex
%
% Hier werden allgemeine Einstellungen festgelegt.

%Packages sind Erweiterungen des eigentlichen Programms und erleichtern Abläufe, Darstellungen, Formatierungen usw.
%Im Folgenden sind die Packages installiert, welche ich zum erstellen einer Bachelorarbeit einmal benötigt wurden.
%Es kann beliebig erweitert/reduziert werden.

%%%%%%%%%%%%%%%%%%%%%%%%%%%%%%%%%%%%%%%%%%%%%%%%%%%%%%%%%
%%%%%%%%%%%%%%%%%%%%%%%%%%%%%%%%%%%%%%%%%%%%%%%%%%%%%%%%%

% T1-Schriftkodierung: nötig, damit Wörter mit Umlauten korrekt getrennt werden
% und der Text im PDF durchsuch-/kopierbar ist (u.a. für die Plagiatsprüfung)
\usepackage[T1]{fontenc}
% Times New Roman Alternative, pdflatex-kompatibel
\usepackage{newtxtext}

% Paket für Times Schrift
%\usepackage{mathptmx}

\usepackage[ngerman]{babel}
% Paket für Deutsche Sprache (Übersetzungen von Chapter zu Kapitel,
% richtige Umlaute, richtige % Silbentrennung)
% siehe auch http://de.wikipedia.org/wiki/Babel-System

\usepackage{titling}
%Hilft im Dokument auf den deklarierten Autor und Titel zuzugreifen

\usepackage{csquotes}
% Zitierweise wird in Berichtinfos.tex über \harvardtrue / \harvardfalse gewählt.
% Anforderung: einheitlicher Stil, Seitenangabe obligatorisch; Online-PDFs mit URL oder DOI
% (daher doi=true).
\ifharvard
    % Kurzzitierweise nach Harvard in Fußnoten, wie im Anforderungsdokument empfohlen:
    %   Fußnote:              siehe Weber (2020): S. 8.
    %   Literaturverzeichnis: Weber, Max (2020): Titel. Ort: Verlag.
    %   Sortierung:           alphabetisch nach Verfasser, dann Jahr, dann Titel (nyt)
    % Verwendung im Text:     \footcite[Vgl.][8]{schluessel}   bzw. \footcite[8\psq]{schluessel} (=S. 8 f.)
    \usepackage[style=authoryear,backend=biber,sorting=nyt,doi=true,isbn=false,
        maxcitenames=2,maxbibnames=3,dashed=false,uniquename=false,uniquelist=false]{biblatex}
    \DeclareDelimFormat{nameyeardelim}{\addspace}
    \DeclareDelimFormat{postnotedelim}{\addcolon\space}
    \DeclareDelimFormat[bib]{nameyeardelim}{\addspace}
    \DeclareDelimFormat[bib]{labelnamepunct}{\addcolon\space}
    \renewbibmacro*{cite:labeldate+extradate}{%
        \iffieldundef{labelyear}{}%
            {\printtext[bibhyperref]{\printtext[parens]{\printlabeldateextra}}}}
\else
    % Numerische Zitierweise (IEEE-ähnlich), Seitenangabe trotzdem Pflicht: \cite[8]{schluessel}
    \usepackage[style=numeric-comp,backend=biber,doi=true,isbn=false,maxnames=3,sorting=none]{biblatex}
\fi
% Internetquellen: Webadresse in eckigen Klammern plus Datum des letzten Zugriffs
\DeclareFieldFormat{url}{\mkbibbrackets{\url{#1}}}
\DefineBibliographyStrings{ngerman}{urlseen = {Zugriff am}}
\addbibresource{Referenzen.bib}
%Regelt die Erstellung des Literaturverzeichnisses sowie die Zitation im gesamten Bericht
%Greift gleichzeitig auf die Referenzen.bib genannte BibTex datei zu, in welcher die Quellen aufgeführt sind.

\usepackage{url}
\urlstyle{same}
%erzeugt schönere URLs (vorallem in den Quellen wichtig), in der Textschrift statt Schreibmaschinenschrift

\usepackage{float}
%Hilft beim Anordnen von Bildern, Grafiken und Tabellen

%\usepackage{abstract}
%ermöglicht das einfügen eines Abstract (Nicht relevant fürs Praktikum)

\usepackage{booktabs}
%ermöglicht das erstellen von schöneren Tabellen

\usepackage{array}
\usepackage{longtable}
%Tabellen, die über mehrere Seiten umbrechen dürfen (z.B. KI-Verzeichnis)

\usepackage{subfigure}
%lässt den Benutzer mehrere Grafiken/Bilder in einer Abbildung darstellen

\usepackage{setspace}
% Paket um den Zeilenabstand zu ändern
\onehalfspacing
% Zeilenabstand auf 1,5-fach setzen
% Zeilenabstand 1,5 Zeilen für den Fließtext (Anforderung).
% Hinweis: \onehalfspacing = 1,5 x Schriftgröße (18 pt bei 12 pt). Wer exakt Words
% "1,5 Zeilen" bei Times New Roman 12 (ca. 20,7 pt) nachbilden will: \setstretch{1.43}

\usepackage{pdfpages}
%ermöglicht das einbinden von eigenständigen PDF Dokumenten

\usepackage{siunitx}
\sisetup{locale = DE}
%vereinfacht den benötigten Syntax zum Benutzen von SI Einheiten und Mathematischen Gleichungen mit dem deutschen Standart

\usepackage{tikz}
%Freies Anordnen von Bildern und Text in Relation zueinander, sowie Möglichkeiten Bilder zu bearbeiten
%Verfügt auserdem über die Möglichkeit Grafiken und Diagramme in Latex zu erstellen

\usepackage[format=plain,labelfont=bf]{caption}
\captionsetup[table]{position=above}
%verschönert die Abbildungs- und Tabellenbeschriftungen. Für Tabellen werden die Beschriftungen über der Tabelle angezeigt

\usepackage{icomma}
%Für den korrekten Abstand der Kommas in Dezimalzahlen und in Gleichungen

\usepackage{amsmath}
%elementares Paket für Gleichungen. Verschönert diese und erleichtert die Benutzung von Gleichungen.

\usepackage{graphicx}
%mehr Optionen beim Einbinden und anordnen von Bildern und Grafiken

%%% HJP 31/3/25
\newlength\dist
\setlength\dist{-\topskip}
\let\chapback\chapterheadstartvskip
\def\setchapspace#1{\setlength\dist{#1\dist}}
\def\lesschapspace{\renewcommand*\chapterheadstartvskip{\vspace*{\dist}}}
\def\resetchapspace{\let\chapterheadstartvskip\chapback}
%%%

\usepackage[left=3cm,right=2cm,top=3cm,bottom=2cm,footskip=1cm]{geometry}
%Regelt die Seitenabstände gemäß Anforderung: links und oben je 3 cm, rechts und unten je 2 cm
%(setzt einseitiges Layout voraus, siehe Option oneside in main.tex)

\usepackage{enumitem}

\usepackage{etoolbox}
%Ermöglicht das Entfernen des Seitenumbruchs vor \chapter (siehe Abschnitt unten)

\usepackage[hidelinks]{hyperref}
% Generische Links

\usepackage[acronym, toc]{glossaries}

\renewcommand*{\glspostdescription}{}
%Entfernt die Nummerierung im Abkürzungsverzeichnis

\usepackage{listings}
\definecolor{codegreen}{rgb}{0,0.6,0}
\definecolor{codegray}{rgb}{0.5,0.5,0.5}
\definecolor{codepurple}{rgb}{0.58,0,0.82}
\definecolor{backcolour}{rgb}{0.95,0.95,0.92}
\lstdefinestyle{myStyle}{
    backgroundcolor=\color{backcolour},
    commentstyle=\color{codegreen},
    keywordstyle=\color{magenta},
    numberstyle=\tiny\color{codegray},
    stringstyle=\color{codepurple},
    basicstyle=\ttfamily\footnotesize,
    breakatwhitespace=false,
    breaklines=true,
    keepspaces=true,
    numbers=left,
    numbersep=5pt,
    showspaces=false,
    showstringspaces=false,
    showtabs=false,
    tabsize=2,
}

\usepackage{todonotes}
% Kann mit disable komplett ausgeschaltet werden
%\usepackage[disable]{todonotes}
\setuptodonotes{fancyline}
\todototoc

\flushbottom

\iffalse % don’t use titlesec.sty with KOMA script classes
\usepackage{titlesec}
%% \titleformat{⟨command⟩}[⟨shape⟩]{⟨format⟩}{⟨label⟩}{⟨sep⟩}{⟨before⟩}[⟨after⟩]
\titleformat{name=\chapter}[display]
{\raggedleft\chaptertitlename\ {\textcolor{gray}{\fontsize{60}    {65}\selectfont\thechapter}}}
{20pt}{}[]
\fi

%%%%%%%%%%%%%%%%%%%%%%%%%%%%%%%%%%%%% ToC DEPTH LEVEL
\setcounter{secnumdepth}{3} % number subsubsection
\setcounter{tocdepth}{3} % list subsubsection


\addtokomafont{chapterprefix}{\raggedleft}

%%%%%%%%%%%%%%%%%%%%%%%%%%%%%%%%%%%%% KAPITEL: BEGINN AUF NEUER SEITE?
% KOMA-Script beginnt jedes Kapitel auf einer neuen Seite (\clearpage, bei "openright" zusätzlich \cleardoublepage).
%
% Schalter:
%   \kapitelneuseitefalse  Die Kapitel laufen fort: mehrere Kapitel koennen
%                          auf derselben Seite stehen. Der Seitenumbruch vor
%                          \chapter wird dazu entfernt - Nummerierung, Format
%                          und Inhaltsverzeichnis der Kapitel bleiben
%                          unveraendert.
%   \kapitelneuseitetrue   Jedes Kapitel beginnt auf einer neuen Seite
%
% Die Kapitel werden in main.tex ueber \kapitel{...} eingebunden: im Standard
% mit \include (eigene .aux-Datei je Kapitel, Seitenumbruch am Dateiende), im
% Fliesstext mit \input (kein erzwungener Seitenumbruch am Dateiende).
\newif\ifkapitelneuseite
\kapitelneuseitetrue
% \kapitelneuseitefalse

\ifkapitelneuseite
  % Standard: Jedes Kapitel beginnt auf einer neuen Seite - nichts zu tun.
\else
  % Fliesstext: Der Seitenumbruch vor \chapter wird entfernt.
  \makeatletter

  \patchcmd{\scr@startchapter}
    {\if@openright\cleardoublepage\else\clearpage\fi}
    {}
    {\relax}
    {%
      \PackageError{WAB-Vorlage}{%
        Der Seitenumbruch vor \string\chapter konnte nicht entfernt werden%
      }{%
        Die interne Definition von \string\scr@startchapter (KOMA-Script) hat%
        sich geaendert. Bitte die Vorlage entsprechend anpassen.%
      }%
    }

  \renewcommand*{\scr@chapter@beforeskip}{1.5\baselineskip}%

  \providecommand*{\scr@chapter@afterindent}[3]{\@afterindentfalse}%

  \renewcommand*{\chapterheadstartvskip}{\vspace{\@tempskipa}}%
  \makeatother
\fi

% Kapitel-Dateien einbinden: \kapitel{Kapitel/1Einleitung}
%   Fliesstext: \input (kein erzwungener Seitenumbruch am Dateiende, damit
%               mehrere Kapitel auf derselben Seite stehen koennen)
%   Standard:   \include (eigene .aux-Datei je Kapitel, Seitenumbruch am
%               Dateiende, daher beginnt jedes Kapitel auf einer neuen Seite)
\ifkapitelneuseite
  \newcommand*{\kapitel}[1]{\include{#1}}
\else
  \newcommand*{\kapitel}[1]{\input{#1}}
\fi

%%%%%%%%%%%%%%%%%%%%%%%%%%%%%%%%%%%%% PREFACE AND BODY STYLING
\def\preface{
    \pagenumbering{roman}
}

\def\body{
    %\cleardoublepage
    \ifkapitelneuseite
      \lesschapspace
    \fi
    \pagenumbering{arabic}
    \pagestyle{headings}
}

\def\abstract{
    \begin{center}{
        \large\bfseries Kurzzusammenfassung}
    \end{center}
    \normalsize
    \normalfont
}
\def\endabstract{
  \par
}

\newenvironment{acknowledgements}{
   \cleardoublepage
    \begin{center}{
        \large \bf Danksagung}
    \end{center}
    \normalsize
    }{\cleardoublepage}
\def\endacknowledgements{
  \par
}

%%%%%%%%%%%%%%%%%%%%%%%%%%%%%%%%%%%%%%%%%%%%%%%%%%%%%%%%%
%%%%%%%%%%%%%%%%%%%%%%%%%%%%%%%%%%%%%%%%%%%%%%%%%%%%%%%%%

%%%%%%%%%%%%%%%%%%%%%%%%%%%%%%%%%%%%% ZITATE AUS KI-QUELLEN
% Anforderung: Wörtliche KI-Zitate *in Asterisken* statt Anführungszeichen,
% Fußnote verweist auf den Eintrag im KI-Verzeichnis (z.B. "ChatGPT 1.").
% Paraphrasen / KI als Ausgangspunkt: Fußnote "Vgl. Elicit 1."
% Übernahme von Strukturen: Fußnote an der Kapitelüberschrift.
%   \kizitat{Text der KI}{ChatGPT 1}
%   \kivgl{Elicit 1}
\newcommand{\kizitat}[2]{*#1*\footnote{#2.}}
\newcommand{\kivgl}[1]{\footnote{Vgl. #1.}}
